\inicializarSeccion{Problemática}

\tab De acuerdo con la Organización Mundial de la Salud \parencite{OMS2025}, la malaria es de las enfermedades infecciosas de mayor impacto a nivel mundial. De acuerdo con la OMS, en 2024 se reportaron 282 millones de casos de malaria y 610 mil muertes en 80 países diferentes, lo que demuestra que la enfermedad sigue siendo un problema grave de salud pública aún con los avances en prevención y tratamiento.

\tab La región más vulnerable por causa de la malaria es la región africana, la cual cuenta con alrededor del 95\% de casos y muertes causadas por la malaria. Este problema afecta principalmente a niños menores de 5 años quienes tienen el 75\% de las muertes de malaria en la región. A su vez, niños, niñas, mujeres embarazadas, viajeros y personas con VIH y sida son quienes están en mayor peligro de infectarse y tener complicaciones graves.

\tab Las soluciones para lograr erradicar el problema son variadas, como prevenir picaduras de mosquitos y tomar medicinas (quimioprofilaxis), para reducir picaduras de mosquitos se pueden utilizar mosquiteros al dormir en zonas donde esté presente la malaria, usar repelentes de mosquitos y usar ropa protectora. Otra de las mejores soluciones es el uso de las vacunas contra malaria RTS,S/AS01 y R21/Matrix-M las cuales reducen de forma significativa la malaria. Esta ya se está dando alrededor de África donde junto a los otros métodos de prevención se espera salvar decenas de miles de vidas cada año. La OMS espera con estas medidas poder reducir los casos de malaria y su tasa de mortalidad un 90\% para 2030, junto con eliminarla en al menos 35 países por completo.

\begin{figure}[H]
    \centering
    \includegraphics[width=0.8\textwidth]{\nat{parte1/mapaMalaria}}
    \caption{Mapa de casos de malaria en el mundo \parencite{WHO2019Malaria}}
    \label{fig:mapaMalaria}
\end{figure}